% GENERATED FILE. Edit canonical lessons, then run npm run generate:books.
% canonical-source-hash: fnv1a64:45c16e1fb28423ee
% canonical-lessons: ZH-C337-gongxi, ZH-C337-zhu-wish, ZH-W337-zhu-wish, ZH-C337-xinyi, ZH-C337-haoyun, ZH-C337-xuyuan

\chapter{Best Wishes}
\label{ch:zh-best-wishes}
\hlchaptermodality{337}

\begin{chapteropening}
\textbf{By the end of this chapter:} \emph{I can congratulate someone, wish them well and good luck, give a token of my regard and make a wish.}

Complete the last lesson of chapter 337: I can congratulate someone, wish them well and good luck, give a token of my regard and make a wish.
\end{chapteropening}

\section[gōngxǐ]{gōngxǐ — to congratulate}

\label{lesson:ZH-C337-gongxi}



\begin{quote}
Before the new one: say the Mandarin for to give a present, then the Mandarin for to shake hands with someone.
\end{quote}

\subsection*{You'll want to know: gōngxǐ}
\emph{gōngxǐ} — \textquotedblleft{}to congratulate\textquotedblright{}.

\begin{quote}
\emph{gōngxǐ nǐ zhǎodào le xīn gōngzuò} — Congratulations on finding a new job.
\end{quote}

Its characters come later; for now it is a word for the ear.

\begin{cousinweb}
\emph{gōngxǐ} is built from \emph{gōng} and \zh{喜} 'joy', both inherited from Middle Chinese, \emph{kjowng} and \emph{xiX} (Guangyun, in Baxter's transcription), which glosses the first 'respectful' and the second 'joy'.
\end{cousinweb}

\begin{grammarlens}[title={What you've built}]
One more word for telling what happened.
\end{grammarlens}

\subsection*{Your turn}
\begin{itemize}
\raggedright
  \item \emph{Say it:} \emph{gōngxǐ}
  \item \emph{Say it:} \emph{gōngxǐ}, once more
  \item \emph{Say it:} \emph{gōngxǐ nǐ zhǎodào le xīn gōngzuò} — Congratulations on finding a new job.
  \item \emph{Recall:} say the Mandarin for to give a present, then the Mandarin for to shake hands with someone, then say \emph{gōngxǐ} again
\end{itemize}

\begin{tcolorbox}[breakable,colback=teal!4,colframe=teal!35!black,title={Before you move on}]
What does gōngxǐ mean? (\textquotedblleft{}To congratulate\textquotedblright{}.) Say it once more.
\end{tcolorbox}
\section[\texorpdfstring{\zh{祝} (zhù)}{祝 (zhù)}]{\zh{祝} (zhù) — to wish someone well}

\label{lesson:ZH-C337-zhu-wish}



\begin{quote}
Before the new one: say the Mandarin for to shake hands with someone, then the Mandarin for to congratulate someone.
\end{quote}

\subsection*{You'll want to know: \zh{祝}}
\textbf{\zh{祝}} — \emph{zhù} — \textquotedblleft{}to wish someone well\textquotedblright{}.

\begin{quote}
\textbf{\zh{祝你生日快乐}} \emph{zhù nǐ shēngrì kuàilè} — Happy birthday to you!
\end{quote}

The new character, in pieces: \zh{祝} — nine strokes: on the left a spirit piece of four strokes, a dot, a stroke that runs right and turns down to the left, a vertical and a dot; on the right \zh{口} over \zh{儿}. The next lesson puts it in your hand.

\begin{cousinweb}
\textbf{\zh{祝}} (\emph{zhù}) is inherited from Middle Chinese \emph{tsyuwH} (Guangyun, in Baxter's transcription), which glosses it, quoting an older dictionary, as the words of praise spoken by the master of a sacrifice. The dictionary also lists \emph{tsyuwk}, 'one who calls on the spirits'; the falling tone of zhù matches the first.
\end{cousinweb}

\begin{grammarlens}[title={What you've built}]
One more word for telling what happened.
\end{grammarlens}

\subsection*{Your turn}
\begin{itemize}
\raggedright
  \item \emph{Say it:} \emph{zhù}
  \item \emph{Say it:} \emph{zhù}, once more
  \item \emph{Say it:} \emph{zhù nǐ shēngrì kuàilè} — Happy birthday to you!
  \item \emph{Recall:} say the Mandarin for to shake hands with someone, then the Mandarin for to congratulate someone, then say \emph{zhù} again
\end{itemize}

\begin{tcolorbox}[breakable,colback=teal!4,colframe=teal!35!black,title={Before you move on}]
What does \zh{祝} mean? (\textquotedblleft{}To wish someone well\textquotedblright{}.) Say it once more.
\end{tcolorbox}
\section[\texorpdfstring{\zh{祝} (zhù)}{祝 (zhù)}]{\zh{祝} — the spirit piece, then \zh{口} over \zh{儿}}

\label{lesson:ZH-W337-zhu-wish}



\begin{quote}
The word first, then the character inside it.

\begin{itemize}
\raggedright
  \item \emph{Recall:} say \emph{zhù} — \textbf{R1}, one lesson back, before it decays
\end{itemize}

\begin{quote}
\zh{口} \zh{儿}
\end{quote}

\emph{Read it:} each one aloud, with nothing said first — \emph{kǒu}, \emph{ér}
\end{quote}

\subsection*{Script — the spirit piece, then \zh{口} over \zh{儿}}
\hlblockfigure{\detokenize{figures/ZH-W337-zhu-wish-filmstrip.pdf}}{How \zh{祝} is written, stroke by stroke}

\begin{quote}
\zh{祝}
\end{quote}

\textbf{Two pieces, left before right.}

\textbf{The left piece, the spirit piece, four strokes.} A dot. A bar that runs right and turns down to the left. A vertical under it. Then a short dot at its right.

\textbf{The right piece, five strokes.} First \zh{口}: a short vertical, a bar that turns down, the bar that closes it. Under it \zh{儿}: a stroke falling to the left, then a vertical that bends right along the foot and hooks up.

4 + 5 = nine.

\subsection*{Your turn}
\begin{itemize}
\raggedright
  \item \emph{Write it:} \zh{祝}, all nine, with the model in front of you
  \item \emph{Cover:} write it again from memory, and count the pieces, not the strokes
\end{itemize}

\begin{tcolorbox}[breakable,colback=teal!4,colframe=teal!35!black,title={Before you move on}]
How many strokes in \zh{祝}? (\textbf{Nine}.) What is on the right? (\textbf{\zh{口} over \zh{儿}}.) How does the spirit piece end? (\textbf{With a short dot}.)

Source: \href{https://raw.githubusercontent.com/chanind/hanzi-writer-data/68d10a4b21150cae5e1ebbd223eed289cf32d90c/data/\%E7\%A5\%9D.json}{Hanzi Writer Data \zh{祝}, pinned stroke paths and medians}.
\end{tcolorbox}
\section[\texorpdfstring{\zh{心意} (xīnyì)}{心意 (xīnyì)}]{\zh{心意} (xīnyì) — kind thoughts; a token of regard}

\label{lesson:ZH-C337-xinyi}



\begin{quote}
Before the new one: say the Mandarin for to congratulate someone, then the Mandarin for to wish someone well.
\end{quote}

\subsection*{You'll want to know: \zh{心意}}
\textbf{\zh{心意}} — \emph{xīnyì} — \textquotedblleft{}kind thoughts; a token of regard\textquotedblright{}.

\begin{quote}
\textbf{\zh{我们送你一个红包这是我们的一点儿心意}} \emph{wǒmen sòng nǐ yí ge hóngbāo zhè shì wǒmen de yìdiǎnr xīnyì} — We're giving you a red envelope, a small token of our regard.
\end{quote}

\begin{cousinweb}
\textbf{\zh{心意}} (\emph{xīnyì}) is built from \zh{心} 'the heart' and \zh{意} 'intent', both inherited from Middle Chinese, \emph{sim} and \emph{'iH} (Guangyun, in Baxter's transcription), which glosses the second 'the will'.
\end{cousinweb}

\begin{grammarlens}[title={What you've built}]
One more word for telling what happened.
\end{grammarlens}

\subsection*{Your turn}
\begin{itemize}
\raggedright
  \item \emph{Say it:} \emph{xīnyì}
  \item \emph{Say it:} \emph{xīnyì}, once more
  \item \emph{Say it:} \emph{wǒmen sòng nǐ yí ge hóngbāo zhè shì wǒmen de yìdiǎnr xīnyì} — We're giving you a red envelope, a small token of our regard.
  \item \emph{Recall:} say the Mandarin for to congratulate someone, then the Mandarin for to wish someone well, then say \emph{xīnyì} again
\end{itemize}

\begin{tcolorbox}[breakable,colback=teal!4,colframe=teal!35!black,title={Before you move on}]
What does \zh{心意} mean? (\textquotedblleft{}Kind thoughts; a token of regard\textquotedblright{}.) Say it once more.
\end{tcolorbox}
\section[hǎoyùn]{hǎoyùn — good luck}

\label{lesson:ZH-C337-haoyun}



\begin{quote}
Before the new one: say the Mandarin for to wish someone well, then the Mandarin for a token of regard.
\end{quote}

\subsection*{You'll want to know: hǎoyùn}
\emph{hǎoyùn} — \textquotedblleft{}good luck\textquotedblright{}.

\begin{quote}
\emph{zǒu yǐqián dàjiā dōu zhù wǒ hǎoyùn} — Before I left, everyone wished me good luck.
\end{quote}

Its characters come later; for now it is a word for the ear.

\begin{cousinweb}
\emph{hǎoyùn} is built from \zh{好} 'good' and \emph{yùn}, both inherited from Middle Chinese, \emph{xawX} and \emph{hjunH} (Guangyun, in Baxter's transcription), which glosses the second 'to move' and 'to carry'. The sense 'luck' is not given there.
\end{cousinweb}

\begin{grammarlens}[title={What you've built}]
One more word for telling what happened.
\end{grammarlens}

\subsection*{Your turn}
\begin{itemize}
\raggedright
  \item \emph{Say it:} \emph{hǎoyùn}
  \item \emph{Say it:} \emph{hǎoyùn}, once more
  \item \emph{Say it:} \emph{zǒu yǐqián dàjiā dōu zhù wǒ hǎoyùn} — Before I left, everyone wished me good luck.
  \item \emph{Recall:} say the Mandarin for to wish someone well, then the Mandarin for a token of regard, then say \emph{hǎoyùn} again
\end{itemize}

\begin{tcolorbox}[breakable,colback=teal!4,colframe=teal!35!black,title={Before you move on}]
What does hǎoyùn mean? (\textquotedblleft{}Good luck\textquotedblright{}.) Say it once more.
\end{tcolorbox}
\section[xǔyuàn]{xǔyuàn — to make a wish}

\label{lesson:ZH-C337-xuyuan}



\begin{quote}
Before the new one: say the Mandarin for a token of regard, then the Mandarin for good luck.
\end{quote}

\subsection*{You'll want to know: xǔyuàn}
\emph{xǔyuàn} — \textquotedblleft{}to make a wish\textquotedblright{}.

\begin{quote}
\emph{shēngrì nà tiān tā xiān xǔyuàn ránhòu chī dàngāo} — On his birthday he made a wish first and then ate the cake.
\end{quote}

Its characters come later; for now it is a word for the ear.

\begin{cousinweb}
\emph{xǔyuàn} is built from \emph{xǔ} and \emph{yuàn}, both inherited from Middle Chinese, \emph{xjoX} and \emph{ngjwonH} (Guangyun, in Baxter's transcription), which glosses the first 'to permit' and 'to grant', and the second 'to desire'.
\end{cousinweb}

\begin{grammarlens}[title={What you've built}]
One more word for telling what happened.
\end{grammarlens}

\subsection*{Your turn}
\begin{itemize}
\raggedright
  \item \emph{Say it:} \emph{xǔyuàn}
  \item \emph{Say it:} \emph{xǔyuàn}, once more
  \item \emph{Say it:} \emph{shēngrì nà tiān tā xiān xǔyuàn ránhòu chī dàngāo} — On his birthday he made a wish first and then ate the cake.
  \item \emph{Recall:} say the Mandarin for a token of regard, then the Mandarin for good luck, then say \emph{xǔyuàn} again
\end{itemize}

\begin{tcolorbox}[breakable,colback=teal!4,colframe=teal!35!black,title={Before you move on}]
What does xǔyuàn mean? (\textquotedblleft{}To make a wish\textquotedblright{}.) Say it once more.
\end{tcolorbox}
